\section{A checked reason to skip separation}\label{sec:screening}

A support oracle should not run merely because a block exists. One
inexpensive certificate can rule out cuts with sufficiently large
normalized violation at the current query. This is a local statement
about a chosen normalization and tolerance, not a prediction of total
solver time.

Let $g(x)=(x,F(x))\in\R^r$ denote the feature vector used by a block,
let $\bar z\in\R^r$ be the current lifted query, and choose positive
coordinate scales $s_1,\ldots,s_r$. Suppose $x^1,\ldots,x^N\in D$
are verified parameter samples and $\lambda_i\ge0$,
$\sum_i\lambda_i=1$. Put $z^*=\sum_i\lambda_i g(x^i)\in H_D$.
With exact graph values, define
\[
\rho_1=\sum_j\frac{|\bar z_j-z_j^*|}{s_j},\qquad
\rho_\infty=\max_j\frac{|\bar z_j-z_j^*|}{s_j}.
\]
The implementation can instead have verified enclosures
$g_j(x^i)\in[\ell_{ij},u_{ij}]$. Define
$L_j=\sum_i\lambda_i\ell_{ij}$, $U_j=\sum_i\lambda_i u_{ij}$, and
\begin{equation}\label{eq:screen}
r_j=\frac{\max\{|\bar z_j-L_j|,|\bar z_j-U_j|\}}{s_j},
\qquad R_1=\sum_jr_j,\quad R_\infty=\max_jr_j.
\end{equation}
All quantities in this expression are rational in the certificate.

\begin{theorem}[Normalized violation screening]\label{thm:screen}
For every valid lower support inequality $c^\top z\ge\beta$ on $H_D$,
\[
\beta-c^\top\bar z\le R_1
\quad\text{if }\max_j s_j|c_j|\le1,
\]
and
\[
\beta-c^\top\bar z\le R_\infty
\quad\text{if }\sum_j s_j|c_j|\le1.
\]
Consequently, $R_1\le\tau$ rules out support cuts with violation
strictly larger than $\tau$ under the first normalization; the
corresponding statement holds for $R_\infty$ under the second.
\end{theorem}
\begin{proof}
The actual convex combination $z^*$ belongs to $H_D$ and satisfies
$L_j\le z_j^*\le U_j$. Validity gives
\[
\beta-c^\top\bar z\le c^\top(z^*-\bar z)
\le\sum_j(s_j|c_j|)\frac{|z_j^*-\bar z_j|}{s_j}.
\]
Bound the last expression by $R_1$ when each first factor is at most
one, or by $R_\infty$ when their sum is at most one. These are the
finite-dimensional H\"older inequalities with an interval upper bound
for each residual.
\end{proof}

Proposed convex weights may come from an ordinary numerical LP. To
obtain a valid witness, interpret its finite coefficients exactly,
replace negative weights by zero, and divide by their positive exact
sum. This preserves a convex combination regardless of the accuracy of
the proposal. An all-nonpositive vector supplies no witness. Each sample
must satisfy the exact box and affine rows, and its original graph value
must be enclosed by the trusted evaluator. A sample feasible only for a
relaxation of the \emph{block domain being certified} is not sufficient
for this screening theorem. If the block deliberately uses a larger
domain than the original model, the conclusion concerns cuts for that
larger-domain hull.

The screen is safe to omit: a large residual, failed LP, rejected sample,
or failed enclosure causes ordinary bounded separation to continue.
For a positive tolerance, a successful screen does not prove exact hull
membership or original nonlinear feasibility. At exactly zero residual,
it proves membership of the query in the chosen block hull, which still
does not imply membership of the query in the original graph or in the
full model's simultaneous hull.

The integration's direction LP bounds each scaled direction coordinate
individually, so the applicable certificate is $R_1$, not $R_\infty$.
The two variants are implemented explicitly in
\pathref{../solver/screening.py}{solver/screening.py}. Sample enclosures
are cached against a fixed evaluator and domain. Serialized replay
re-evaluates the original graph and verifies the proposed mixture, query,
normalization, and threshold; serialized enclosures alone are not trusted.
This theorem is standard convex-combination geometry. Its role is to
give a concrete, checkable skip condition, separate from the heuristic
decision to spend time looking for one.
